% !TeX root = ../深度学习与大语言模型技术原理.tex
% 由主文档通过 \input 载入；请编译主文档。
\setcounter{chapter}{4} % 保留原书章号。
\chapter{Transformer 与现代解码器}
\label{ch:transformer}

\section{从信息选择到缩放点积注意力}

仅靠逐位置 MLP 无法直接混合不同词元的信息。\emph{自注意力（self-attention）}让每个位置根据自身查询，选择其他位置的内容。设 $X\in\mathbb R^{n\times d}$ 按行存储表示，取投影矩阵 $W_Q,W_K\in\mathbb R^{d_h\times d}$、$W_V\in\mathbb R^{d_v\times d}$，定义
\begin{equation}
Q=XW_Q^\top,\quad K=XW_K^\top,\quad V_X=XW_V^\top,\qquad
O=\operatorname{softmax}_{\mathrm{row}}\!\left(\frac{QK^\top}{\sqrt{d_h}}+M\right)V_X.
\label{eq:attention}
\end{equation}
这里用 $V_X$ 表示值矩阵，以区别词表大小 $V$。查询（query）决定需要什么信息，键（key）决定匹配关系，值（value）是被加权汇总的内容。每一行的注意力权重非负且和为1，因而该行输出位于值向量的凸包；后续投影和残差连接不受这一凸包限制。

为什么除以 $\sqrt{d_h}$？假设初始化附近，$q_j,k_j$ 跨坐标独立、均值为零、方差为1，并且两组向量独立，则 $\operatorname{Var}(q^\top k)=d_h$。缩放后方差为1，可减轻维度增大导致 softmax 过度饱和的问题。真实训练中的坐标并不始终满足这些假设，因此这是合理的尺度设计依据，而不是任意模型激活的精确分布结论\cite{attention}。

\section{因果掩码、多头与键值共享}

仅解码器模型（decoder-only model）要求位置 $i$ 只能读取 $j\le i$，因此取 $M_{ij}=0$ 当 $j\le i$，否则取 $-\infty$。实现中还要结合填充掩码与样本边界掩码。对于全被屏蔽的行，softmax 没有合法分布，不能简单用常规指数归一化；应避免产生这种查询或使用有明确定义的内核行为。

\emph{多头注意力（multi-head attention, MHA）}使用 $h_q$ 组查询、键和值投影。取常用的 $d_v=d_h$，把各头输出按列拼接为 $O_{\mathrm{cat}}\in\mathbb R^{n\times h_qd_h}$，再计算 $O_{\mathrm{cat}}W_O^\top$，其中 $W_O\in\mathbb R^{d\times h_qd_h}$。不同头可以学习不同匹配关系，但“一个头对应一个语法规则”不是结构本身保证的事实。\emph{多查询注意力（multi-query attention, MQA）}让所有查询头共享一组键和值；\emph{分组查询注意力（grouped-query attention, GQA）}令 $h_q$ 个查询头分为 $h_{kv}$ 组，每组共享键和值\cite{gqa}。

当 $h_q=h_{kv}$ 时得到 MHA，当 $h_{kv}=1$ 时得到 MQA。若 $h_{kv}$ 整除 $h_q$，每个键值头服务 $h_q/h_{kv}$ 个查询头。共享直接减少键值投影和缓存，但查询头仍各自计算匹配分数，因此不能据此把所有注意力运算量都按同一比例缩小。缓存公式将在式~\eqref{eq:kv-memory} 中给出。

\section{残差、归一化与门控前馈层}

对列向量 $x\in\mathbb R^d$，层归一化（Layer Normalization, LayerNorm）计算
\[
\mu=\frac1d\sum_jx_j,\quad v=\frac1d\sum_j(x_j-\mu)^2,\quad
\operatorname{LN}(x)=g\odot\frac{x-\mu\mathbf1}{\sqrt{v+\varepsilon}}+b.
\]
均方根归一化（Root Mean Square Layer Normalization, RMSNorm）省去减均值，常写为
\[
\operatorname{RMSNorm}(x)=g\odot\frac{x}{\sqrt{d^{-1}\sum_jx_j^2+\varepsilon}}.
\]
二者都按单个词元的特征维归一化，不依赖同批其他样本的均值。RMSNorm 并不使输出均值为零，学习缩放 $g$ 后也不能声称输出均方根恒为1\cite{rmsnorm}。

前置归一化（Pre-Norm）块可写为
\[
u=x+\operatorname{Attn}(\operatorname{Norm}(x)),\qquad
y=u+\operatorname{FFN}(\operatorname{Norm}(u)).
\]
后置归一化（Post-Norm）则对残差求和后的结果归一化。Pre-Norm 的残差路径在求导时保留直接恒等分量，但最终训练稳定性仍依赖深度、初始化和优化设置；不能仅凭顺序判断哪一种在所有任务上更优。

前馈网络（feed-forward network, FFN）通常逐位置执行。门控线性单元（gated linear unit, GLU）的 SwiGLU 变体为
\begin{equation}
\operatorname{FFN}(x)=W_2\bigl(\operatorname{SiLU}(W_1x)\odot W_3x\bigr),
\label{eq:swiglu}
\end{equation}
其中 $W_1,W_3\in\mathbb R^{d_f\times d}$、$W_2\in\mathbb R^{d\times d_f}$。两条输入投影经逐元素门控交互，再映回隐藏维度\cite{swiglu}。忽略偏置，其参数量为 $3dd_f$；与普通两层 FFN 比较时要固定参数或计算预算，不能只固定中间宽度。

\section{旋转位置编码的相对位置结构}

若不输入位置信息，无掩码自注意力对位置置换具有等变性，难以区分词序。\emph{旋转位置编码（rotary position embedding, RoPE）}把每个查询和键的相邻两维看成二维平面，在位置 $m$ 施加角度 $m\omega_j$ 的旋转\cite{rope}。设头维度为偶数，对第 $j$ 个二维块定义
\[
R_j(m)=\begin{pmatrix}\cos(m\omega_j)&-\sin(m\omega_j)\\
\sin(m\omega_j)&\cos(m\omega_j)\end{pmatrix},
\qquad R(m)=\operatorname{diag}(R_1(m),\ldots,R_{d_h/2}(m)).
\]
频率可取 $\omega_j=b^{-2(j-1)/d_h}$，其中 $b>1$ 是基数。旋转后的查询键为 $\widetilde q_m=R(m)q_m$、$\widetilde k_n=R(n)k_n$。

\begin{proposition}
\label{prop:rope-relative}
若查询和键使用同一组频率，则
\[
\widetilde q_m^\top\widetilde k_n=q_m^\top R(n-m)k_n.
\]
\end{proposition}
\begin{proof}
二维旋转满足 $R_j(m)^\top=R_j(-m)$ 及 $R_j(a)R_j(b)=R_j(a+b)$。分块对角乘法因此给出 $R(m)^\top R(n)=R(n-m)$，代入内积即可。
\end{proof}

命题说明显式旋转项依赖相对位移；$q_m,k_n$ 本身仍依赖内容与之前各层的上下文，不能说整个得分只由距离决定。RoPE 保持向量范数，也不天然保证超出训练长度后的性能。推理缓存通常保存已经旋转的键，位置编号必须与训练及缓存偏移一致；长上下文扩展见第~\ref{ch:longcontext} 章。

\section{计算量与一个完整解码器块}

标准稠密注意力的分数矩阵有 $n^2$ 个元素；在固定头数且 $h_qd_h=d$ 时，$QK^\top$ 与权重乘值的运算量为 $O(n^2d)$，投影与 FFN 通常为 $O(nd^2)$ 或 $O(ndd_f)$。直接保存所有注意力矩阵需 $O(h_qn^2)$ 空间；是否必须把它们写入显存，则取决于实现，第~\ref{ch:efficientattention} 章将讨论在线 softmax。

完整模型依次执行词元嵌入、各层归一化与带 RoPE 的因果注意力、残差和 SwiGLU、最终归一化以及词表投影。训练在所有有效位置计算损失，推理只需取最新位置的输出分布。验证小型实现时，应先检查改变未来词元不会影响过去位置的输出，再检查全序列前向与缓存增量前向在相同数值条件下相符。

